\section{Rational relation modules}
\label{sec:modules}

Because \(A_i:M_i^\QQ\to P^\QQ\) is a surjective
\(\QQ\GammaG\)-map, semisimplicity gives
\[
 [K_i]=[M_i^\QQ]-[P^\QQ]
\]
in the rational representation ring.  Exact character inner products
therefore determine \(K_i\) from the block and point permutation
characters.

\begin{proposition}[Complete relation-module table]
\label{prop:relation-table}
The multiplicities of the eight rational types in \(K_1,K_2,K_3\) are as
follows.
\[
\begin{array}{c|rrr|r}
V&K_1&K_2&K_3&\dim_{\QQ}V\\
\midrule
\mathbf1 &343&354&355&1\\
\eps     &341&352&351&1\\
\iot     &314&304&302&2\\
\om      &339&349&353&2\\
\eps\om  &338&347&349&2\\
\iot\om  &314&304&302&4\\
\rh      &1312&1314&1310&16\\
\rh\om   &1305&1304&1306&32\\
\midrule
\multicolumn{1}{r|}{\text{weighted total}}
&66\,674&66\,674&66\,674&
\end{array}
\]
In particular, \(K_1,K_2,K_3\) are pairwise nonisomorphic as
\(\QQ\GammaG\)-modules.
\end{proposition}

\begin{proof}
Subtract the point multiplicities in \cref{prop:point-module} from the
three exact block permutation characters.  The resulting table is the
one displayed.  For example, the two large-type contributions are
\[
 16(1312,1314,1310)\quad\text{for }\rh,
 \qquad
 32(1305,1304,1306)\quad\text{for }\rh\om.
\]
Summing every row with its rational dimension
gives \(66\,674\) in each column, agreeing with the incidence rank.

The three entries in every row are not all equal; more strongly, each
pair of columns differs in at least one multiplicity, and inspection shows
that the three configurations are separated in every rational type.
Uniqueness of semisimple multiplicities proves pairwise
nonisomorphism.
\end{proof}

\begin{remark}[What the table separates]
The trivial multiplicities \(343,354,355\) are the dimensions of the
invariant relation spaces and can also be recovered from block-orbit
counts.  The structural conclusion is the simultaneous decomposition of
the full \(66\,674\)-dimensional kernels.  It should not be confused with
equality of incidence images: all three maps have the same full rational
image \(P^\QQ\), while their kernels are inequivalent.
\end{remark}
